\section{Neural networks for sequences}
\label{sec:bg:sequences}

The stateful models of the thesis read a stream of book events and keep a state per
market.  This section builds them up from their parts: the multilayer perceptron, the
recurrent network and the gated recurrent unit that updates a state, the way such a
network is trained through time, the attention mechanism and the transformer encoder that
read a window of book states, the continuous-time treatment of a state between irregular
events, and the pretrained time-series models that are used without training.

\subsection{The multilayer perceptron}
\label{sec:bg:mlp}

The multilayer perceptron (MLP) is the building block of every network here: a sequence
of \emph{layers}, each a linear map followed by an activation function $\sigma$ applied
element by element, with the last layer linear.  With one hidden layer,
\begin{equation}
  \mathbf{a} = \mathbf{W}_1 \mathbf{x} + \mathbf{b}_1, \qquad
  \mathbf{h} = \sigma(\mathbf{a}), \qquad
  \mathbf{z} = \mathbf{W}_2 \mathbf{h} + \mathbf{b}_2, \qquad
  \hat{\mathbf{p}} = \operatorname{softmax}(\mathbf{z}),
  \label{eq:mlp}
\end{equation}
where $\mathbf{W}_1 \in \mathbb{R}^{d \times n}$ and $\mathbf{W}_2 \in \mathbb{R}^{3
\times d}$ are matrices of weights and $\mathbf{b}_1$, $\mathbf{b}_2$ vectors of biases;
deeper networks stack more layers, $\mathbf{h}^{(l)} = \sigma(\mathbf{W}^{(l)}
\mathbf{h}^{(l-1)} + \mathbf{b}^{(l)})$ with $\mathbf{h}^{(0)} = \mathbf{x}$, and
without the non-linear $\sigma$ the layers would collapse into one linear map.  The
gradient of the loss with respect to the parameters of \cref{eq:mlp} is the backward pass
of \cref{sec:bg:learning:backprop} written out: $\partial \ell / \partial \mathbf{z} =
\hat{\mathbf{p}} - \mathbf{e}_y$, then $\partial \ell / \partial \mathbf{W}_2 =
(\partial \ell / \partial \mathbf{z})\, \mathbf{h}^{\top}$, $\partial \ell / \partial
\mathbf{h} = \mathbf{W}_2^{\top}\, \partial \ell / \partial \mathbf{z}$, $\partial \ell /
\partial \mathbf{a} = \partial \ell / \partial \mathbf{h} \odot \sigma'(\mathbf{a})$ and
$\partial \ell / \partial \mathbf{W}_1 = (\partial \ell / \partial \mathbf{a})\,
\mathbf{x}^{\top}$, each line reusing the one before.  A layer
with $d$ units maps its input to $d$ numbers; the \emph{width} $d$ and the number of
layers, the \emph{depth}, are the two sizes of an MLP, and its parameters are the entries
of the matrices and biases, $\sum_l d_l (d_{l-1} + 1)$ in all.  With one hidden layer of
enough units an MLP can approximate any continuous function on a bounded domain, but the
functions it learns easily are the smooth ones: an MLP is a poor learner of sharp
thresholds on single inputs, which is where decision trees are strong
(\cref{sec:bg:tabular}), and a good learner of smooth interactions between many inputs.
The predictor of the stateful model is an MLP of one hidden layer of 64 units on the
concatenation of its states; the tabular networks of \cref{sec:bg:tabular} are MLPs of
three layers of 512 units.

\subsection{Recurrent networks and the gated recurrent unit}
\label{sec:bg:rnn}

An MLP reads one vector and forgets it.  A \emph{recurrent neural network} reads a
sequence $\mathbf{x}_1, \mathbf{x}_2, \ldots$ one element at a time and keeps a
\emph{state} $\mathbf{h}_t \in \mathbb{R}^{d}$ that summarises what it has read so far,
updated at every step from the previous state and the new input by the same function with
the same parameters,
\begin{equation}
  \mathbf{h}_t = \tanh(\mathbf{W} \mathbf{x}_t + \mathbf{U} \mathbf{h}_{t-1} + \mathbf{b}),
  \label{eq:rnn}
\end{equation}
and the prediction at step $t$ is read from $\mathbf{h}_t$.  The state is the memory of
the network: whatever the prediction at step $t$ needs from the past has to be in
$\mathbf{h}_t$, written there by the updates of the previous steps.  The simple update of
\cref{eq:rnn} has a well-known flaw.  The gradient of a loss at step $t$ with respect to
the state at step $t - k$ is a product of $k$ Jacobians $\partial \mathbf{h}_{s} /
\partial \mathbf{h}_{s-1}$, each of which contains $\mathbf{U}$ and the derivative of the
$\tanh$; a product of $k$ matrices shrinks or grows geometrically with $k$, so the
gradient either \emph{vanishes}, and the network cannot learn to use anything more than a
few steps back, or \emph{explodes} \cite{bengio1994vanishing,pascanu2013difficulty}.
Exploding gradients are tamed by \emph{clipping}, rescaling the gradient vector whenever
its norm exceeds a threshold (1.0 in the stateful model); vanishing ones are the reason for
gated units.

The \emph{gated recurrent unit} (GRU) \cite{cho2014gru}, a lighter cousin of the LSTM
\cite{hochreiter1997lstm}, replaces \cref{eq:rnn} with an update in which the network
decides, at every step and for every component of the state, how much of the old state to
keep and how much of a new candidate to write (\cref{fig:bg:gru}):
\begin{align}
  \mathbf{r}_t &= \sigma(\mathbf{W}_r \mathbf{x}_t + \mathbf{U}_r \mathbf{h}_{t-1} + \mathbf{b}_r)
  && \text{(reset gate)}, \notag \\
  \mathbf{z}_t &= \sigma(\mathbf{W}_z \mathbf{x}_t + \mathbf{U}_z \mathbf{h}_{t-1} + \mathbf{b}_z)
  && \text{(update gate)}, \notag \\
  \tilde{\mathbf{h}}_t &= \tanh\bigl(\mathbf{W}_h \mathbf{x}_t + \mathbf{U}_h (\mathbf{r}_t \odot \mathbf{h}_{t-1}) + \mathbf{b}_h\bigr)
  && \text{(candidate)}, \notag \\
  \mathbf{h}_t &= (1 - \mathbf{z}_t) \odot \mathbf{h}_{t-1} + \mathbf{z}_t \odot \tilde{\mathbf{h}}_t
  && \text{(interpolation)},
  \label{eq:gru}
\end{align}
where $\sigma$ is the logistic sigmoid, so that the gates are vectors of numbers in
$(0, 1)$.  A component of the state whose update gate is near zero is copied unchanged
from one step to the next, and the gradient flows back through that copy without
shrinking, which is what lets a GRU keep information over many steps; the reset gate lets
the candidate ignore the old state where that is useful.  A GRU has three times the
parameters of \cref{eq:rnn}, $3 d (n + d + 1)$ for an input of size $n$.  In the thesis a
GRU appears in two roles: as an \emph{encoder} that reads the $K$ book states of a window
one after the other and returns its last state (\cref{sec:models:encoders}), and as a
\emph{cell} that performs one update of a market's state from one input vector
(\cref{sec:models:ctdgnn}); in the second role it is written $\operatorname{GRUCell}(\mathbf{x},
\mathbf{h})$ and \cref{eq:gru} is applied once.

\begin{figure}[t]
\centering
\resizebox{\textwidth}{!}{\begin{tikzpicture}[>=Latex, x=1cm, y=1cm,
  box/.style={draw=vizInk, rounded corners=2pt, align=center, font=\footnotesize, inner sep=3pt, minimum height=9mm, text width=44mm, fill=white},
  input/.style={box, dashed, text width=26mm}, op/.style={box, fill=vizB!12}, state/.style={box, fill=vizA!12, text width=26mm},
  arr/.style={->, draw=vizInk, line width=0.8pt}, lab/.style={font=\scriptsize, text=vizInk}]
\node[input] (x)    at (0,1.3)    {input $\mathbf{x}_t$};
\node[state] (hp)   at (0,-1.3)   {previous state $\mathbf{h}_{t-1}$};
\node[op]    (r)    at (5.6,2.4)  {reset gate\\$\mathbf{r}_t = \sigma(\mathbf{W}_r \mathbf{x}_t + \mathbf{U}_r \mathbf{h}_{t-1} + \mathbf{b}_r)$};
\node[op]    (z)    at (5.6,0.0)  {update gate\\$\mathbf{z}_t = \sigma(\mathbf{W}_z \mathbf{x}_t + \mathbf{U}_z \mathbf{h}_{t-1} + \mathbf{b}_z)$};
\node[op]    (cand) at (5.6,-2.4) {candidate\\$\tilde{\mathbf{h}}_t = \tanh(\mathbf{W}_h \mathbf{x}_t + \mathbf{U}_h (\mathbf{r}_t \odot \mathbf{h}_{t-1}) + \mathbf{b}_h)$};
\node[op]    (mix)  at (11.4,0.0) {interpolation\\$\mathbf{h}_t = (1 - \mathbf{z}_t) \odot \mathbf{h}_{t-1} + \mathbf{z}_t \odot \tilde{\mathbf{h}}_t$};
\node[state] (h)    at (15.6,0.0) {new state $\mathbf{h}_t$};
\draw[arr] (x.east) -- (r.west); \draw[arr] (x.east) -- (z.west); \draw[arr] (x.east) -- (cand.west);
\draw[arr] (hp.east) -- (r.west); \draw[arr] (hp.east) -- (z.west); \draw[arr] (hp.east) -- (cand.west);
\draw[arr] (r.east) -- ++(0.5,0) |- node[lab, pos=0.25, right] {$\mathbf{r}_t$} (cand.east);
\draw[arr] (z.east) -- node[lab, above] {$\mathbf{z}_t$} (mix.west);
\draw[arr] (cand.east) -| node[lab, pos=0.25, above] {$\tilde{\mathbf{h}}_t$} (mix.south);
\draw[arr] (hp.south) |- ++(0,-2.4) -| node[lab, pos=0.25, below] {$\mathbf{h}_{t-1}$} (mix.south east);
\draw[arr] (mix) -- (h);
\end{tikzpicture}
}
\caption{The gated recurrent unit.  From the input $\mathbf{x}_t$ and the previous state
$\mathbf{h}_{t-1}$ the two gates and the candidate are computed in parallel; the new state
is an interpolation between the old state and the candidate, component by component,
controlled by the update gate.}
\label{fig:bg:gru}
\end{figure}

\subsection{Training through time}
\label{sec:bg:tbptt}

A recurrent network is trained like any other: the loss is summed over the steps and
backpropagation computes its gradient.  Because the state at step $t$ depends on the
parameters through every previous step, the computation is first \emph{unrolled}, that
is, written as one long chain of copies of the same update, one per step, and the backward
pass runs through the whole chain, summing at each copy the gradient with respect to the
shared parameters (\emph{backpropagation through time}).  On a stream that is days long,
millions of steps, the chain cannot be kept in memory, and the gradient would in any case
vanish beyond a few hundred steps.  \emph{Truncated} backpropagation through time
\cite{williams1990tbptt} cuts the chain into windows of a fixed number of steps
(\cref{fig:bg:tbptt}): the forward pass runs through the whole stream and carries the state
across windows, but the backward pass runs only inside the current window, and the state at
the start of the window is treated as a constant (\emph{detached}), so that no gradient
flows into the previous windows.  The stateful model of the thesis uses windows of five
blocks and accumulates the gradients of twenty windows before an optimiser step; what the
model can learn about dependencies longer than five blocks it learns through the state
itself, which is why the state has to be a good memory on its own, and why the thesis adds
to it the slow memories of \cref{sec:bg:ctime}.

\begin{figure}[t]
\centering
\resizebox{\textwidth}{!}{\begin{tikzpicture}[>=Latex, x=1cm, y=1cm,
  st/.style={draw=vizInk, rounded corners=2pt, align=center, font=\footnotesize, inner sep=3pt, minimum height=8mm, minimum width=14mm, fill=vizA!12},
  inp/.style={draw=vizInk, dashed, rounded corners=2pt, align=center, font=\scriptsize, inner sep=2pt, minimum height=6mm, minimum width=14mm, fill=white},
  outp/.style={draw=vizInk, rounded corners=2pt, align=center, font=\scriptsize, inner sep=2pt, minimum height=6mm, minimum width=14mm, fill=vizC!15},
  arr/.style={->, draw=vizInk, line width=0.8pt},
  grad/.style={->, draw=vizB, line width=0.8pt, dashed},
  cut/.style={draw=vizB, line width=1.2pt},
  lab/.style={font=\scriptsize, text=vizInk, align=center}]
\foreach \i/\t in {0/t-4, 1/t-3, 2/t-2, 3/t-1, 4/t} {
  \node[st]  (h\i) at (2.2*\i, 0)    {$\mathbf{h}_{\t}$};
  \node[inp]  (x\i) at (2.2*\i, 1.4)  {$\mathbf{W}_{\t}$};
  \node[outp] (y\i) at (2.2*\i, -1.4) {$\ell_{\t}$};
  \draw[arr] (x\i) -- (h\i);
  \draw[arr] (h\i) -- (y\i);
}
\node[st] (hprev) at (-2.2, 0) {$\mathbf{h}_{t-5}$};
\draw[arr] (hprev) -- (h0);
\foreach \i/\j in {0/1, 1/2, 2/3, 3/4} { \draw[arr] (h\i) -- (h\j); }
% gradients flow backwards inside the window
\foreach \i/\j in {1/0, 2/1, 3/2, 4/3} { \draw[grad] (h\i.north west) ++(0.1,0.05) to[bend right=25] ([xshift=-0.1cm, yshift=0.05cm]h\j.north east); }
\foreach \i in {0,1,2,3,4} { \draw[grad] (y\i.north) ++(-0.25,0) -- ++(0,0.55); }
% detach at the window boundary
\draw[cut] (-1.1,-1.9) -- (-1.1,2.0);
\node[lab, text=vizB, anchor=south] at (-1.1,2.05) {detach: no gradient\\into the previous window};
\node[lab, anchor=north] at (4.4,-2.0) {one truncation window of five blocks: the losses $\ell_{t-4}, \ldots, \ell_t$ are summed,\\the gradient (dashed) flows back through the five states and stops at the boundary};
\end{tikzpicture}
}
\caption{Truncated backpropagation through time.  The state is carried forward through
the whole stream, but the gradient of the summed losses of a window flows back only through
the states of that window and stops at its boundary, where the incoming state is treated as
a constant.}
\label{fig:bg:tbptt}
\end{figure}

This training regime differs from that of a tabular model in a way that matters for the
results.  A tabular model sees its training rows in random order, several times, in
mini-batches that mix markets and hours; a stateful model sees the stream once, in time
order, and the mini-batch of a step is whatever markets happen to be active in that block.
The gradient estimates of the second regime are far more correlated from one step to the
next, so the optimiser takes a noisier path, and the model has seen every row once when
training ends.  \Cref{sec:stateful:regime} measures what this costs on the same table.

\subsection{Attention and the transformer encoder}
\label{sec:bg:attention}

A recurrent network reads a window of $K$ book states in order and compresses them step by
step; \emph{attention} reads all of them at once and lets every element look at every
other.  Given $n$ vectors, the \emph{tokens} $\mathbf{X} \in \mathbb{R}^{n \times d}$
(the rows), three learned linear maps produce for every token a query
$\mathbf{q}$, a key $\mathbf{k}$ and a value $\mathbf{v}$, and the output for token $i$ is a
weighted average of the values of all tokens, the weights being a softmax of how well its
query matches their keys:
\begin{equation}
  \operatorname{Attn}(\mathbf{X}) = \operatorname{softmax}\!\Bigl(\frac{\mathbf{Q}\mathbf{K}^{\top}}{\sqrt{d_k}}\Bigr)\mathbf{V},
  \qquad
  \mathbf{Q} = \mathbf{X}\mathbf{W}_Q,\ \mathbf{K} = \mathbf{X}\mathbf{W}_K,\ \mathbf{V} = \mathbf{X}\mathbf{W}_V,
  \label{eq:attention}
\end{equation}
where the softmax is taken row by row and $\sqrt{d_k}$ keeps the dot products on a
scale at which the softmax is not saturated \cite{vaswani2017attention}.  Row $i$ of the
result is a mixture of the values of the tokens that are relevant to token $i$, with
weights that the network learns; the operation does not depend on the order of the tokens,
and when the order matters it is added as a \emph{positional encoding} to the inputs.
Several attentions with different weight matrices, the \emph{heads}, run in parallel and
their outputs are concatenated, so that different heads can attend to different things.  A
\emph{transformer encoder layer} (\cref{fig:bg:transformer}) is a multi-head attention
followed by a small MLP applied to each token, each of the two wrapped in a \emph{residual
connection} (the input is added to the output, so that the layer learns a correction
rather than a replacement, and the gradient has a direct path back) and a layer
normalisation.  Stacking such layers gives the transformer, the architecture of large
language models and of the time-series foundation models of \cref{sec:bg:tsfm}.

\begin{figure}[t]
\centering
\begin{tikzpicture}[>=Latex, x=1cm, y=1cm, node distance=5mm,
  box/.style={draw=vizInk, rounded corners=2pt, align=center, font=\footnotesize, inner sep=4pt, minimum height=9mm, text width=66mm, fill=white},
  input/.style={box, dashed}, op/.style={box, fill=vizB!12}, outbox/.style={box, fill=vizC!15},
  arr/.style={->, draw=vizInk, line width=0.8pt}, lab/.style={font=\scriptsize, text=vizInk, align=left}]
\node[input] (in)  {$n$ tokens, each a vector of $d$ numbers\\$\mathbf{X} \in \mathbb{R}^{n \times d}$ (here: one token per book channel)};
\node[op, below=of in] (qkv) {three linear maps\\$\mathbf{Q} = \mathbf{X}\mathbf{W}_Q,\ \mathbf{K} = \mathbf{X}\mathbf{W}_K,\ \mathbf{V} = \mathbf{X}\mathbf{W}_V$};
\node[op, below=of qkv] (att) {scaled dot-product attention, one head\\$\operatorname{Attn}(\mathbf{X}) = \operatorname{softmax}\!\bigl(\mathbf{Q}\mathbf{K}^{\top}/\sqrt{d_k}\bigr)\mathbf{V}$\\(rows: for each token, a weighted average of all tokens)};
\node[op, below=of att] (mh) {$h$ heads in parallel, concatenated and mixed\\$\operatorname{MultiHead} = [\operatorname{Attn}_1, \ldots, \operatorname{Attn}_h]\,\mathbf{W}_O$};
\node[op, below=of mh] (an1) {residual connection and layer normalisation\\$\mathbf{X}' = \operatorname{LN}(\mathbf{X} + \operatorname{MultiHead}(\mathbf{X}))$};
\node[op, below=of an1] (ff) {position-wise feed-forward network\\$\operatorname{FF}(\mathbf{x}) = \mathbf{W}_2\, \sigma(\mathbf{W}_1 \mathbf{x} + \mathbf{b}_1) + \mathbf{b}_2$, applied to each token};
\node[op, below=of ff] (an2) {residual connection and layer normalisation\\$\mathbf{X}'' = \operatorname{LN}(\mathbf{X}' + \operatorname{FF}(\mathbf{X}'))$};
\node[outbox, below=of an2] (out) {$n$ tokens of $d$ numbers, each now informed by the others\\$\mathbf{X}'' \in \mathbb{R}^{n \times d}$};
\draw[arr] (in) -- (qkv); \draw[arr] (qkv) -- (att); \draw[arr] (att) -- (mh); \draw[arr] (mh) -- (an1);
\draw[arr] (an1) -- (ff); \draw[arr] (ff) -- (an2); \draw[arr] (an2) -- (out);
\draw[arr] (in.east) -- ++(0.5,0) |- (an1.east);
\draw[arr] (an1.east) -- ++(0.5,0) |- (an2.east);
\end{tikzpicture}

\caption{One transformer encoder layer.  Every token produces a query, a key and a value;
attention mixes the values of all tokens with weights given by the match between queries
and keys; a feed-forward MLP is then applied to every token; residual connections and
layer normalisation wrap both.}
\label{fig:bg:transformer}
\end{figure}

In the thesis the tokens are not the time steps of a series but its \emph{channels}: the
time-first encoder of \cref{sec:models:encoders} maps each of the 55 book quantities, read
as a series of $K$ values, to one token with a linear map shared by all channels, adds a
learned embedding that says which channel the token is, and lets one attention layer mix
the 55 tokens.  This is the channel-independence pattern of \cref{sec:bg:lob}: time is
read first, within each channel, by a map that is the same for all of them, and the
channels are mixed afterwards.  With 55 tokens of 32 numbers, four heads and one layer the
encoder is small (about 40 thousand parameters); its cost is in the attention matrix,
$55 \times 55$ per market per block.

\subsection{Continuous time: decay and multi-rate memories}
\label{sec:bg:ctime}

The updates above assume a sequence of steps with no notion of how much time passes between
them, whereas a market's book events arrive at irregular times and a block may contain an
update for a market whose previous event was a second or a day ago.  Two devices give the
state a sense of time.  The first is a \emph{decay}: between two events the state is
multiplied by $e^{-r \Delta t}$, where $\Delta t$ is the elapsed time and $r > 0$ a rate,
so that an old state fades towards zero and the network reads a faded state as ``little is
known''.  This is the solution of the differential equation $\mathrm{d}\mathbf{h} /
\mathrm{d}t = -r \mathbf{h}$, the simplest continuous-time dynamics, and it has the useful
property that decaying over $a$ and then over $b$ seconds is the same as decaying once over
$a + b$, so the decay can be applied lazily, at a node's next event, without visiting the
node at every block.  GRU-D \cite{che2018grud} introduced a learned decay of this kind for
irregularly sampled medical series; in the stateful model the rate $r$ is a single learned
scalar (\cref{sec:models:decay}), and it converges to a half-life of seventeen seconds.
The second device gives the model several memories at several speeds.  A single state
updated every two seconds and decaying in seconds cannot remember the morning; the
\emph{clockwork} RNN \cite{koutnik2014clockwork} partitions the state into modules that
update at different rates, so that a slow module keeps a summary of the distant past while
a fast one tracks the present.  The hierarchical memory of \cref{sec:models:memory} is a
continuous-time version of this idea: besides the per-block state, each market has a
minute state and an hour state, each a GRU cell that reads the state one level below once
per period, and the predictor reads the three side by side.

\subsection{Time-series foundation models}
\label{sec:bg:tsfm}

A \emph{foundation model} for time series is a transformer pretrained once on a very large
and varied corpus of series and then used on new series \emph{zero-shot}, that is, without
any training on them: given a context window of past values it outputs a forecast of the
next values.  Chronos \cite{ansari2024chronos} quantised the values of a series into
tokens and trained a language model on them; Chronos-2 \cite{ansari2025chronos2}, the model
used here, replaces the tokens with \emph{patches}, short segments of the series mapped to
vectors, and is trained with a \emph{quantile loss}: for each quantile level $\tau$ among
nine (0.1 to 0.9) and each future step, the output $\hat{q}$ is trained to minimise
$\max\bigl(\tau (y - \hat{q}),\, (\tau - 1)(y - \hat{q})\bigr)$, the pinball loss, whose
minimiser is the $\tau$-quantile of the distribution of $y$; the forecast is therefore
not one number but a band of nine quantiles, from which a direction, a ``no change'' zone
and a width can be read (\cref{sec:models:zeroshot}).  Every context window is standardised
before it enters the model, which is what makes the model indifferent to the scale of a
series and, on a one-tick process, what makes it read noise as signal
(\cref{sec:external:tsfm}).  The evidence of such models in finance is sobering: zero-shot
forecasts of returns are close to a random walk, with directional accuracy near 50\%, and
the value found so far comes from a learned correction on top of a frozen model
\cite{rahimikia2025revisiting,dewage2026hybrid}.  This is the setting tested in
\cref{sec:external:tsfm}.
